% !TEX program = context
% !TEX root = somerville-naryunion-sizes.tex

% Allow the specimen to use the reduced big-operator export.
\definefontfeature[mathematics][mathematics][mathexpansion=no]
\definefontfeature[math][math][mathexpansion=no]
\definefontfeature[math-text][math-text][mathexpansion=no]
\definefontfeature[math-script][math-script][mathexpansion=no]
\definefontfeature[math-scriptscript][math-scriptscript][mathexpansion=no]

\input type-imp-somerville.mkxl

% Exported user coordinates for the Agate SemiCondensed Upright master.
\defineSomervilleAxisTypeface{somerville-agate-proof}
  {wdth=88,wght=900,MGHT=0,opsz=5,STYA=5,STYB=5,ARLN=200,ARHD=100,INSL=0}
  {wdth=88,wght=900,opsz=5}
\setupbodyfont[somerville-agate-proof,18pt]
\setuppapersize[A4]
\setuplayout[backspace=24mm,topspace=22mm,width=middle,height=middle,header=0pt,footer=0pt]
\setuppagenumbering[location=none]

\definefont[NaryUnionLabel][Serif at 10pt]
\definefont[NaryUnionTitle][Serif at 25pt]
\definefont[NaryUnionIndexFont][Serif at 5pt]
\definecolor[NaryUnionIndexColor][s=0.75]
\unexpanded\def\NaryUnionNumber#1{%
  {\NaryUnionIndexFont\color[NaryUnionIndexColor]{#1}}%
}

% ConTeXt numbers the base glyph as 1; the sample indices are 0–20.
\unexpanded\def\NaryUnionSample#1{%
  \vbox{%
    \offinterlineskip
    \hbox to \NaryUnionColumnWidth{\hss
      \m{\textstyle\Umathchar 0 0
        \getmathvariant\numexpr#1+1\relax
        "22C3\relax}%
      \hss}%
    \vskip 0.7mm
    \hbox to \NaryUnionColumnWidth{\hss\NaryUnionNumber{#1}\hss}%
  }%
}

\unexpanded\def\NaryUnionColumn#1#2{%
  \bTD
    \vbox{%
      \offinterlineskip
      \def\NaryUnionColumnWidth{#1}%
      \def\NaryUnionGap{}%
      \processcommalist[#2]\NaryUnionStackItem
    }%
  \eTD
}

\unexpanded\def\NaryUnionStackItem#1{%
  \NaryUnionGap
  \NaryUnionSample{#1}%
  \def\NaryUnionGap{\vskip 0.7mm}%
}

\starttext
\midaligned{\NaryUnionTitle N-ary union}
\blank[small]
\midaligned{\NaryUnionLabel Somerville Agate SemiCondensed}
\blank[big]

\startalignment[middle]
\bTABLE[frame=off,align={middle,lohi},offset=0pt,distance=0.5mm]
  \bTR
    \NaryUnionColumn{15mm}{0,13}
    \NaryUnionColumn{21mm}{20,1}
    \NaryUnionColumn{15mm}{12,1}
    \NaryUnionColumn{20mm}{2,18}
    \NaryUnionColumn{15mm}{2,11}
    \NaryUnionColumn{19mm}{16,0}
  \eTR
\eTABLE
\blank[1mm]
\bTABLE[frame=off,align={middle,lohi},offset=0pt,distance=0.5mm]
  \bTR
    \NaryUnionColumn{18mm}{1,15}
    \NaryUnionColumn{13mm}{10,3}
    \NaryUnionColumn{19mm}{17,0}
    \NaryUnionColumn{12mm}{4,9}
    \NaryUnionColumn{20mm}{2,19}
    \NaryUnionColumn{11mm}{8,5}
    \NaryUnionColumn{17mm}{14,1}
    \NaryUnionColumn{11mm}{6,7}
  \eTR
\eTABLE
\stopalignment

\startbuffer[naryunion-prose]
\startalignment[flushleft,verytolerant]
Lorem ipsum dolor sit amet, consectetur adipiscing elit, sed do eiusmod tempor
incididunt ut labore et dolore magna aliqua. Ut enim ad minim veniam, quis
nostrud exercitation ullamco laboris nisi ut aliquip ex ea commodo consequat.
\stopalignment
\stopbuffer

\blank[big]
\doloopoverlist{0,1,2,3,4,5,8,10,15,20}{%
  \edef\NaryUnionSize{\recursestring}%
  \getbuffer[naryunion-prose]\par
  \startformula
    \dorecurse{\ifnum\NaryUnionSize>10 3\else 5\fi}{%
      \ifnum\recurselevel>1 \thinspace\fi
      {\rm ABC} \Umathchar 0 0 \getmathvariant\numexpr\NaryUnionSize+1\relax
        "22C3\relax
    }%
    \quad\hbox{\NaryUnionNumber{\NaryUnionSize}}%
  \stopformula
}

\stoptext
